\section{Very large prime incidence on the actual boundary arms}
\label{vlg-section}

Let \(n>324\) be prime, \(B=n^4\), \(r=\lceil\sqrt{6n}\rceil\),
and \(L=\log(8B)\). Retain the actual roots
\[
 I_B=\{k\in\mathbb Z:B\leq k<2B\},\qquad
 a_k=3k,\quad w_k=a_k+\zeta,
\]
\[
 Q_k=a_k^2+a_k+1,\quad \alpha_k=w_k/\bar w_k,
 \qquad \zeta^2=\zeta-1.
\]
Fix a subset \(I\subseteq I_B\) whose actual ratios are
multiplicatively independent. The established private-norm set
or any fixed color from Theorem~\ref{ncs-colors} supplies such a set.
No independence assertion for an arbitrary subset of the whole block
is made.

Write \(w_k^n=A_k+C_k\zeta\), and set
\[
 b_{k,1}=|A_k|,\quad b_{k,2}=|C_k|,\quad
 b_{k,3}=|A_k+C_k|,\qquad
 T_k=\prod_{j=1}^3b_{k,j},\quad
 c_k=\max_j b_{k,j},\quad t_k=\log c_k.
\]
These powers are primitive and unramified, so the arm sizes are
positive and pairwise coprime. A zero boundary would make \(w_k^n\)
a unit times a rational integer, contradicting the nonzero oriented
split-prime valuations of \(w_k/\bar w_k\).
The actual bounds \(9B^2\leq Q_k<36B^2\) give
\begin{equation}\label{vlg-height}
 n\log(3B)\leq t_k<nL.
\end{equation}
Indeed, \(c_k\geq Q_k^{n/2}\), while
\[
 c_k\leq\frac2{\sqrt3}Q_k^{n/2}
       <\frac2{\sqrt3}(6B)^n<(8B)^n.
\]
The norm and each arm are coprime: reducing the norm modulo a prime
dividing \(A_k\), \(C_k\), or \(A_k+C_k\) leaves the square of a
nonzero coordinate.

\begin{theorem}[Two endpoints at very large primes]\label{vlg-endpoints}
For every rational prime \(q\) satisfying
\begin{equation}\label{vlg-cutoff}
 \log q\geq \frac r2 L,
\end{equation}
the set \(E_q=\{k\in I:v_q(T_k)\geq4\}\) has at most two elements.
\end{theorem}
\begin{proof}
The cutoff gives \(q>8B>n\) and \(q^4\geq(8B)^{2r}\).
Every hit supplies a unit norm and unit conjugate denominator
modulo \(q^4\). The identity
\[
 (A+C\zeta)^3-(A+C\bar\zeta)^3
   =3AC(A+C)(\zeta-\bar\zeta)
\]
therefore implies \(\alpha_k^{3n}=1\pmod{q^4}\).
Since \(q\nmid3n\), the established split/inert simple-lifting count
bounds this norm-one torsion group by \(3n\).

The coordinate calculation in the proof of Lemma~\ref{um-rigidity}
gives, for two actual products of \(r\) roots, a determinant bounded by
\[
 |D|\leq8r\,7^{2r-1}B^{2r-1}<(8B)^{2r}.
\]
The strict inequality is uniform in \(r\): after division, with
\(x=49/64\), its upper bound is
\[
 \frac8{7B}r x^r
 <\frac8{7B}\frac1{1-x}
 =\frac{512}{105B}<1.
\]
Equal ratio products modulo \(q^4\) thus give \(q^4\mid D\), hence
\(D=0\) and exact ratio equality in \(\mathbb Q(\zeta)\).
Independence recovers every multiplicity. Consequently the multisets
of length \(r\) from \(E_q\) inject into the torsion group, giving
\[
 \binom{|E_q|+r-1}{r}\leq3n
\]
when \(E_q\) is nonempty. If \(|E_q|\geq3\), its left side is at least
\[
 \binom{r+2}{2}=\frac{(r+1)(r+2)}2>\frac{r^2}2\geq3n,
\]
a contradiction. Empty hit sets cause no difficulty.
\end{proof}

\begin{theorem}[Actual simultaneous packets and arm degrees]
\label{vlg-packets}
Let \(V=I\times\{1,2,3\}\). Take as labels precisely the primes
satisfying \eqref{vlg-cutoff} for which \(E_q\ne\varnothing\), and put
\[
 e(q)=\{(k,j):k\in E_q,\ q\mid b_{k,j}\}.
\]
Each label has one or two endpoints, with distinct roots in the
two-endpoint case. For an incident endpoint \(v=(k,j)\), write
\(\epsilon_v(q)=v_q(b_{k,j})\geq4\).
For every subset \(U\subseteq V\) and every subset \(P\) of the labels,
\begin{equation}\label{vlg-divisibility}
 \prod_{q\in P}q^{\,\sum_{v\in U\cap e(q)}\epsilon_v(q)}
       \ \bigm|\ \prod_{v\in U}b_v.
\end{equation}
Writing \(d_U(q)=|U\cap e(q)|\), this implies
\begin{equation}\label{vlg-weighted}
 4\sum_q d_U(q)\log q
 \leq\sum_q\sum_{v\in U\cap e(q)}\epsilon_v(q)\log q
 \leq\sum_{v\in U}\log b_v.
\end{equation}
With \(d=\lfloor n/(2r)\rfloor\), every arm vertex has degree at most
\(d\), and every root has at most \(3d\) incident prime labels.
\end{theorem}
\begin{proof}
The arm sizes are pairwise coprime, so each hit root belongs to
exactly one endpoint at this prime. Theorem~\ref{vlg-endpoints}
then gives the endpoint count. The label set is finite, since its
members divide the fixed nonzero integer \(\prod_{k\in I}T_k\).
Different prime labels may have the same endpoint set, and singleton
edges are allowed.

For each prime \(q\), the exponent on the left of
\eqref{vlg-divisibility} is at most its valuation in the integer on
the right. Applying this simultaneously to the distinct primes
proves the divisibility. Taking logarithms proves
\eqref{vlg-weighted}. Primes omitted from the graph, including
depths below four at other vertices, only add nonnegative mass on
the right.

If \(m\) labels meet a fixed arm \(v=(k,j)\), then
\[
 2rLm\leq\sum_{q\ni v}4\log q
       \leq\log b_v\leq t_k<nL.
\]
Thus \(m<n/(2r)\). Since \(n\) is odd and \(2r\) is even, this
ratio is not an integer, giving \(m\leq d\). Summing over the
three different arms gives the root-degree bound.
\end{proof}

The total number of labels is at most \(3d|I|\), by summing
incidences. If every label had two endpoints, it would be at most
\(3d|I|/2\); singleton labels have not been excluded. In particular,
the available cardinality bound is \(O(|I|\sqrt n)\), and does not
provide an \(o(B)\) estimate.

\begin{theorem}[A joint matching decomposition]\label{vlg-matchings}
The entire actual prime-labelled graph on arm vertices is a disjoint
union of at most \(2d-1\) matchings. Within each matching, different
prime labels have disjoint endpoint sets. They may use different
arms of the same root.
\end{theorem}
\begin{proof}
Here \(d\geq1\), since
\(2r\leq2\sqrt{6n}+2<n\) for \(n>324\).
Process the finite labels in increasing order. At an edge with two
endpoints, at most \(2(d-1)\) earlier labels meet either endpoint:
each endpoint has total degree at most \(d\), including the current
edge. A singleton edge has at most \(d-1\) earlier neighbors.
Multiple edges only cause overcounting in these upper bounds.
Thus one of \(2d-1\) colors is absent from all earlier adjacent
labels. Assign the first available color. Its color classes are the
claimed matchings. If \(I\) is empty, there are no labels to color.
This is a direct finite construction, without a multigraph-coloring
or matroid-partition assumption.
\end{proof}

This statement applies to all of the actual labels simultaneously.
It does not bound the total weight of a matching or resolve the
independence of roots outside the established domain.

\begin{proposition}[The remaining weighted budget]\label{vlg-net}
For every subset \(U\) of the actual arm vertices,
\begin{equation}\label{vlg-net-equation}
 \sum_q\sum_{v\in U\cap e(q)}(\epsilon_v(q)-3)\log q
 \leq\sum_{v\in U}\log b_v-3\sum_q d_U(q)\log q.
\end{equation}
For \(U=V\), its left side is precisely the positive excess in the
region \eqref{vlg-cutoff}.
\end{proposition}
\begin{proof}
Subtract \(3\sum_q d_U(q)\log q\) from the second inequality of
\eqref{vlg-weighted}. Each incidence has depth at least four.
Pairwise coprimality allocates the entire depth of \(T_k\) to its
unique arm, proving the last assertion.
\end{proof}

This retains the graph's radical saving, but does not account for
the negative contributions at depths one and two outside this
positive graph. The estimates do not make its right side
\(o(\sum_{k\in I}t_k)\).

For a precise limitation of the finite inequalities alone, consider
an abstract weighted ledger for sufficiently large \(n\). Give each
of \(N\) formal roots three arms and
\(m=\lfloor n/(4r)\rfloor\) private labels on each arm. Give every
label depth four and logarithmic weight \(w=(r/2)L\). All edges are
singletons, and the full mass on an arm is
\[
 4mw=2rmL\leq nL/2.
\]
Its degree is at most \(d\), and every moment support at a fixed
label has cardinality one. Yet its positive excess per root is
\[
 3mw\sim \frac38 nL.
\]
Thus the degree, moment-cardinality and strict arm-cap inequalities
alone permit a linear positive cost. These formal labels and
weights have not been realized by distinct rational primes or by
Eisenstein powers. This is not a counterexample to any actual
root-family statement, to signed compensation, or to \(abc\).
The stronger simultaneous integer structure in
\eqref{vlg-divisibility} remains a possible source of improvement.

The outstanding interfaces are actual private-label counts and
weights, interactions between different primes, the omitted signed
credits, and the roots outside the proved independent domain.
None is declared automatically satisfied or impossible.
All results in this section are ordinary proofs. No new formal
Lean statement, prime-distribution certificate or density experiment
is claimed.
